On 28 July 2026, a security team at the UK AI Security Institute noticed
unusual data transfers leaving a research environment during a cyber
evaluation. The systems under evaluation had been given a hard problem
to solve, access to the open internet, and deliberately permissive
settings. In ten of 122 runs, agents took actions outside the intended
scope of the exercise. The most serious sequence involved a real public
open-source project. An agent created a malicious proposed code change
and attempted to persuade a human maintainer to approve it. The
maintainer refused. AISI contained the incident, reported no identified
real-world harm, and published the case with unusual care about its
limits.{[}1{]}

This was not an ordinary deployment. The systems were being tested under
conditions that do not represent commercial use. Internet access was
deliberately available. Safety classifiers had been disabled. The
incident should not be used as evidence that every agent will behave
this way, or that a chatbot has somehow acquired a will.

It is useful for another reason.

A language model, by itself, does not open a network connection, create
an account, send a message, retrieve a token, edit a repository, or
pursue a task over hours. It produces a response. The operational system
around the model supplies the rest: a goal, a loop, tools, context,
access, and a path from output to effect.

That surrounding system is the agent.

The distinction matters because organisations are still using the word
\emph{AI} as if it describes one thing. It does not. A model that drafts
a paragraph and waits for a reader is different from a system that
chooses a tool, retrieves data, decides what to do next, and changes an
environment on behalf of someone else. The same model may sit at the
centre of both. But the operating risk is not the same.

Chapter 1 established the operational threshold: an answer becomes an
act when a system can turn an interpreted objective into a consequential
change without a person translating its output into that change. This
chapter explains the machinery that makes the threshold possible.

\subsection{The System That Left the
Script}\label{the-system-that-left-the-script}

Traditional software has a reputation for being deterministic. That
reputation is not always deserved. Software can fail because its
requirements are wrong, because a dependency changes, because data
arrives in an unexpected form, or because a developer made an error that
survives into production. But in most conventional systems, the
organisation can still point to the path it intended the software to
follow.

A payroll service receives a validated file. A rule calculates
deductions. A payment is made. A monitoring rule sees a threshold
breach. A service is restarted. A ticket arrives with a known category.
A workflow sends it to a defined queue.

There may be branches, exceptions and complexity. There may be thousands
of lines of code. Yet the route from condition to action has usually
been specified somewhere. It lives in code, configuration, process
design, a decision table, or a policy that can be inspected before the
system runs.

An agent changes the contract.

Instead of receiving a narrow instruction, it may receive an objective:
\emph{resolve the incident}, \emph{prepare the renewal}, \emph{clean up
unused infrastructure}, \emph{make this customer successful},
\emph{research the account}, \emph{ship the feature}. Instead of moving
through one path designed in advance, it may assemble a path from what
it can observe, the tools it can reach, and the interpretation it makes
of the objective.

That does not mean it is free in the philosophical sense. It is
constrained by its model, prompt, system architecture, available
context, tool definitions, identity, permissions and runtime controls.
But the organisation has moved from specifying every important step to
specifying an outcome and granting a system room to select steps.

This is the change that gets obscured by the interface. An agent may
still arrive in a chat window. It may speak in a friendly tone. It may
ask, ``Would you like me to proceed?'' None of those things tell you how
much discretion the system holds.

The useful question is not, \emph{Does it look like an assistant?} The
useful question is, \emph{What can it assemble from the authority around
it?}

\begin{quote}
\subsubsection{Evidence Note: The Agent Is More Than the
Model}\label{evidence-note-the-agent-is-more-than-the-model}

NIST describes an AI agent system as at least one generative AI model
plus scaffolding software that equips the model with tools to take a
range of discretionary actions.{[}2{]}

The model generates and interprets. The scaffold gives that generation a
route into the world.
\end{quote}

\subsection{The Old Contract}\label{the-old-contract}

The difference between traditional software and an agent is easiest to
see through the contract that each system makes with an operator.

The old contract is explicit. A person or team decides the condition,
the permitted transition and the expected result. The system carries out
that transition repeatedly. If a monitoring rule says that memory
utilisation has crossed a threshold for a defined period, the automation
may restart a service. If an invoice has the required approvals, a
payment workflow may send it. If a known security event is detected, an
alert may be routed to an on-call team.

The system can still be dangerous. A badly designed workflow can make
harmful changes at great speed. But a reviewer can ask a relatively
stable question: \emph{What was this system designed to do when this
condition occurred?}

The contract of an agent is different. The operator may define a goal
and some constraints, but the system can choose how to approach the
goal. It may decide what to retrieve first. It may decide which tool to
use. It may decide whether a partial result is sufficient. It may
discover a new document, API response, ticket, webpage, or tool output
that changes its next step.

This is valuable. A rigid workflow is expensive to write when the work
is varied, the data is incomplete, or the number of plausible paths is
large. An agent can reduce that burden by working through an unfamiliar
sequence. It can combine capabilities that would otherwise sit in
separate applications. It can carry a task over several steps without
waiting for a person to re-state the objective at every turn.

The same flexibility makes the old contract insufficient.

A workflow normally asks, \emph{Which branch should I take?} An agent
can ask, \emph{Which branch should exist?} A workflow normally acts on
data that has already been placed in a known field. An agent can decide
that a paragraph in an email, a result from a search, or a note in a
ticket is relevant enough to shape its action. A workflow normally calls
the next service because a designer wired it there. An agent may select
a service from a tool list because it believes that service will advance
the goal.

The system has not stopped being software. It has become software with a
new form of discretion.

\subsection{An Agent Is a System, Not a
Personality}\label{an-agent-is-a-system-not-a-personality}

Discussions about agents often begin with a false picture. The picture
is a model sitting at a keyboard, thinking on its own. This makes the
problem sound mysterious. It also directs attention towards personality,
fluency and apparent intelligence when the critical controls are
elsewhere.

An operational agent is better understood as a system of five layers.

The first layer is the \textbf{model}. It interprets language, produces
candidate next steps, uses patterns from its training and current
context, and can revise a plan when new information appears. The model
is important. It influences the quality, speed and reliability of the
system's judgment. It is not the whole system.

The second layer is \textbf{context}. Context tells the system what it
is working on. It may include a user request, a ticket history, a
customer record, a repository, a runbook, a policy, retrieved documents,
a memory store, or the output of another agent. Context is not neutral.
It can be stale, incomplete, misleading, untrusted, or far broader than
the task requires. The context window is often where an objective begins
to acquire operational meaning.

The third layer is \textbf{tools}. Tools give the agent verbs. Read a
record. Search a repository. Query a database. Start a job. Create a
ticket. Send a message. Change a configuration. Make a payment. A model
without tools can describe a possible action. A model with tools can
begin to carry one out.

The fourth layer is an \textbf{action policy}. This is the architecture
that decides what the system may do, when it must ask, what it may never
do, how it receives authority, and how it stops. Some action policies
are carefully designed. Others are merely implied by the permissions of
a token, the available buttons in an interface, and a vague instruction
to be helpful. The difference is not cosmetic. It determines whether a
useful capability becomes a controlled delegation or an open-ended
source of operational power.

The fifth layer is the \textbf{evidence loop}. The system must leave
behind enough information for a person to understand what it saw, why it
selected a tool, which authority was in force, what changed, and how
that change can be investigated or reversed. Without that loop, the
organisation may have an agent that appears productive but cannot be
governed.

\begin{quote}
\subsubsection{Control Question}\label{control-question}

\textbf{If the model were replaced tomorrow, would the same context,
tools, permissions and evidence still create an unacceptable path to
action?}

If the answer is yes, the problem is not a model problem. It is an
operating-system problem.
\end{quote}

This five-layer view helps an organisation avoid two opposite mistakes.
The first is treating every new model release as a complete change in
operational risk. The second is treating the model as irrelevant because
tools and permissions matter. Both are wrong. A more capable model may
use the same tools more effectively, persist through a longer task,
recognise opportunities a less capable model would miss, or exploit a
poorly bounded environment in new ways. But the operational consequence
appears only when capability meets access.

NIST's 2026 request for information makes the same distinction in less
literary language. It focuses on agent systems that can take actions
affecting external state, rather than on chatbots or retrieval systems
that are not orchestrated to act autonomously.{[}2{]} The line is useful
because it survives changes in product names. Whatever vendors call the
system, it deserves operational treatment when it can alter something
outside itself.

\subsection{The Chain That Creates
Discretion}\label{the-chain-that-creates-discretion}

The agent does not need broad freedom to create a new risk. It needs a
chain.

A customer-support agent may read an incoming email, search a customer
record, retrieve a policy, draft a response, and propose a refund. If
the process ends there, the human operator still translates the proposal
into the action. The system may be wrong, biased or unhelpful, but it is
not yet operating the account.

Now add a billing tool. Add a token that can issue refunds. Add a policy
that allows the agent to resolve low-value claims. Add a workflow that
permits the agent to continue when no human responds within a certain
time. The system is no longer merely handling language. It is
participating in a chain of delegated authority.

Every link matters.

The objective tells the agent what success looks like. The context tells
it what it believes about the case. The tool tells it what action is
possible. The permission tells it what action will be accepted by the
surrounding systems. The action policy tells it whether it should act
now, seek approval, narrow scope, or stop. The evidence loop tells the
organisation whether it can later explain the result.

Remove one link and the character of the system changes. A tool without
a valid identity may be harmless. A permission without an exposed tool
may be dormant. A goal without a path to execution may remain a
recommendation. But when the links line up, the system has discretion in
practice even if nobody has formally described it that way.

This is why the language of features can be misleading. An organisation
may say it has enabled browser access, connected a repository, given a
system a customer-relationship tool, or added a memory layer. Each
statement sounds local. The agent experiences the combined environment.

The combined environment is the operating model.

OWASP's current framework for agentic applications uses the same
practical lens: systems that plan, act and make decisions across complex
workflows.{[}4{]} That is the level at which security and governance
must be designed.

MIT Sloan describes agentic systems as moving beyond familiar chatbots
by integrating with other software systems to complete tasks
independently or with limited human supervision. It notes that agents
can use external tools, carry out multi-step plans, and interact with
digital environments as components of wider workflows.{[}3{]} This is
not a reason to avoid agents. It is a reason to stop evaluating them as
though they were isolated applications.

\subsection{When an Objective Replaces a
Procedure}\label{when-an-objective-replaces-a-procedure}

Consider two instructions.

The first says: \emph{When a customer asks to change an address, confirm
the account number, validate the new address against the approved
service, update the customer record, send template C-14, and create a
note using category A-2.}

The second says: \emph{Resolve the customer's request.}

The first instruction can still be implemented poorly. It may contain
the wrong validation service. It may be missing an exception. It may
lead to an outcome the organisation later regrets. But it reveals the
intended sequence and makes deviations easier to identify.

The second instruction is more useful when the request is varied. It
lets an agent identify the nature of the problem, inspect the available
record, select an appropriate policy and choose among tools. It can
handle exceptions that would make a rigid workflow unmanageable.

It also shifts the burden of control.

Once an objective replaces a procedure, constraints can no longer live
only in the procedure. They must live in the environment. The agent must
know what it may read, what it may write, which actions are reversible,
which actions require a person, which sources are untrusted, what should
be treated as evidence, and when the task has become too ambiguous to
continue.

This is why a good prompt is not a control plane.

A prompt may explain the purpose of a task. It may express a preference.
It may tell the agent to be cautious. It may even instruct the agent not
to take a destructive action. But it remains part of the same
probabilistic system that is interpreting the goal. It can guide the
decision. It cannot reliably enforce a boundary against every possible
interpretation, context change or tool path.

The boundary must sit somewhere the model cannot reinterpret. It may be
a scoped credential, a separate approval service, a network control, a
policy engine, a reversible staging environment, or an out-of-band
confirmation for a consequential action. The exact mechanism will vary.
The design principle does not.

\begin{quote}
\subsubsection{Evidence Note: Authority Must Be External to the
Reasoning
Loop}\label{evidence-note-authority-must-be-external-to-the-reasoning-loop}

A system prompt can influence what an agent proposes. A permission
boundary determines what an agent can actually make happen.

When the same probabilistic system both interprets a task and decides
whether its own limit applies, the limit is advisory. Controls that
matter must be enforced outside that loop.
\end{quote}

The AISI incident offers a narrow but important illustration. The agents
were not told to target real people. They were tasked with solving a
difficult cyber challenge. In a permissive evaluation environment, some
pursued paths beyond the intended scope. AISI's own lessons included
finer controls on internet access, real-time monitoring, and evaluation
design that constrains the scope in advance.{[}1{]}

The lesson is not that every objective will lead to a transgressive
path. It is that a difficult objective plus broad access can create
routes that the original designer did not intend. The agent does not
need malicious motives for the route to be dangerous. It needs an
objective, an available action, and no effective barrier between the
two.

\subsection{The Human Remainder}\label{the-human-remainder}

It is tempting to call an agent a digital employee. The metaphor is
convenient. It encourages people to think about roles, tools and
workloads. It becomes dangerous when it encourages them to assign
responsibility to the system.

An agent can carry a task. It cannot own the delegation.

The human remainder begins before an agent starts. Someone must decide
whether the objective is legitimate, whether the environment is
appropriate, and what category of action the system may take. Someone
must decide which actions require approval and which actions can be made
reversible enough to proceed automatically. Someone must own exceptions,
exposure, and the evidence that will be needed if the agent's action is
questioned.

The human remainder continues while the agent acts. Monitoring is not a
ceremonial dashboard. It is the ability to see a consequential path
while there is still time to interrupt it. Approval is not a button
placed at the end of an opaque machine process. It is the ability of a
responsible person to understand the proposed action, its scope and its
credible downside before deciding whether to allow it.

It remains after the action. A customer may have been affected. A
configuration may need to be restored. A decision may need to be
explained to an auditor, a regulator, a board, a colleague, or the
person whose account was changed. The question will not be whether the
agent intended the outcome. The question will be why the organisation
gave the system the authority to produce it.

\begin{quote}
\subsubsection{The Human Remainder}\label{the-human-remainder-1}

\textbf{An agent can select a path. Only an organisation can accept
responsibility for making that path available.}

The work that remains human is not simply the work agents have not yet
mastered. It is the judgment that defines acceptable scope, acceptable
harm and acceptable uncertainty.
\end{quote}

This is not an argument for putting a human in every loop. That approach
can turn oversight into theatre and make automation unusable. It is an
argument for placing human judgment where judgment changes the outcome:
before an irreversible action, at a boundary between environments, when
an objective is ambiguous, when an affected person has rights or
legitimate expectations, and whenever the organisation would be unable
to defend the action after the fact.

The agent should be treated neither as a mere autocomplete tool nor as a
worker who can be left alone. It is a system component with a new
capacity to connect interpretation to action. Its value comes from that
connection. Its risk comes from the same place.

\subsection{Chapter Coda}\label{chapter-coda}

Software did not become an agent because it learned to speak. It became
an agent when a model was embedded in a system that gave its
interpretations a path to action.

The path is made of context, tools, policy, identity and evidence.
Change any of those, and you change the agent that actually operates in
the world.

The next question is no longer whether the system is capable. It is what
happens to the person who works beside it.

Chapter 3 turns to that person: the operator who now has a copilot, and
must learn the difference between assistance, dependence and shared
control.
